% Teil II: der durchgerechnete Weg.

\subsection*{Gegeben und gesucht}
Gegeben: Person $m=\qty{70}{kg}$, Boot $M=\qty{140}{kg}$, Bootlänge
$L=\qty{6}{m}$, keine Reibung zwischen Boot und Wasser. Relativgeschwindigkeit
$u=\qty{1}{m/s}$ in Teil (b).
Gesucht: (a) Verschiebung des Bootes beim Gang der Person zum Bug; (b)
Geschwindigkeit des Bootes relativ zum Wasser während und nach dem Gang.

\subsection*{Ansatz}
Ohne äussere horizontale Kräfte bleiben Gesamtimpuls und Schwerpunkt konstant.
Alle Orts- und Geschwindigkeitsangaben deshalb auf das Wasser (Inertialsystem)
beziehen; nur die im Aufgabentext genannte $u$ ist relativ zum Boot.

\subsection*{Rechnung}
\textbf{a) Verschiebung des Bootes.} $x=0$ beim Heck anfangs; Schwerpunkt beider
Massen bleibt am selben Ort wie zu Beginn (alles in Ruhe). Person wandert um
$L$ relativ zum Boot; das Boot verschiebt sich um $s$ entgegengesetzt. Da
innerhalb des Systems nur die Person relativ verschoben wird, gilt
\begin{equation}
  0 = m\,(L-s) + M\,(-s)
  \;\Rightarrow\;
  s = \frac{m}{m+M}\,L
  = \frac{70}{210}\,\qty{6}{m}
  = \qty{2.0}{m}\;.
\end{equation}
($L-s$ ist die Wasser-Verschiebung der Person, $-s$ die des Bootes.)

\textbf{b) Geschwindigkeit des Bootes.} Impulserhaltung bezogen aufs Wasser:
Person hat $u+v$, Boot $v$ (entgegengesetztes Vorzeichen übernehmen):
\begin{equation}
  0 = m(u+v) + Mv
  \;\Rightarrow\;
  v = -\frac{m}{m+M}\,u
  = -\frac{70}{210}\,\qty{1}{m/s}
  = \qty{-0.33}{m/s}\;.
\end{equation}
Das Boot fährt also mit \qty{0.33}{m/s} \emph{entgegen} der Laufrichtung. Bei
$u=0$ (Person bleibt stehen) ist $v=0$: das Boot und die Person ruhen wieder
relativ zum Wasser; die gesamte Verschiebung ist unverändert
$\qty{2}{m}$.

\begin{ergebnis}
(a) Das Boot verschiebt sich um $\dfrac{m}{m+M}\,L=\qty{2.0}{m}$ in Richtung
Bug der Person \emph{entgegen} (also zum Heck). (b) Relativ zum Wasser gilt
\begin{equation*}
  v_\mathrm{Boot}=-\dfrac{m}{m+M}\,u=\qty{-0.33}{m/s},\qquad
  v_\mathrm{Person \text{ zu Beginn}}=u+v=\qty{0.67}{m/s}\;,
\end{equation*}
und sobald die Person stehen bleibt, steht auch das Boot still.
\end{ergebnis}

\begin{kontrolle}
\textbf{Grenzfall:} Für $m\ll M$ (leichte Person, schweres Boot) wird $s
\approx 0$ und $v\approx 0$: das schwere Boot bewegt sich kaum. Für $m=M$
wird $s=L/2$. Plausibel.

\textbf{Einheiten:} $m/(m+M)$ ist dimensionslos, mal Meter ergibt Meter;
Geschwindigkeit $= m/s$. \checkmark

\textbf{Alternative:} über den konstanten Schwerpunkt: derselbe Versatz folgt
aus $\Delta x_S=0$, was die Identität der Ergebnisse bestätigt.
\end{kontrolle}

\begin{merksatz}
In einem System ohne äussere Kräfte genügt der Impulserhaltungssatz (oder
konstanter Schwerpunkt), um zwei Unbekannte zu verknüpfen. Man muss nur
sauber ein gemeinsames Inertialsystem (das Wasser) wählen und relative gegen
absolute Bewegung abgrenzen.
\end{merksatz}
